\section{Discussion}
\label{sec:discussion}

\subsection{A simple, per-scene detector is enough}

The four-detector comparison (Section~\ref{sec:res_fourway}) changes how
the single- versus dual-polarisation question should be posed. Two lessons
follow. First, the choice worth making is not how many polarisations to
use but whether the decision rule is re-estimated for each acquisition:
retraining per scene is what separates the competitive dual SVM from the
fixed-training variant that ranks last, and it also removes any dependence
on a baseline year, settling the objection that a fixed training year is
arbitrary. Second, the second polarisation is not worth its overhead: it
adds no reliable accuracy over a wide range of geometry and climate, and
where it does help, the Sicilian Ancipa case (Section~\ref{sec:res_sicily})
shows the mechanism is real at an extreme but does not generalise.

The cost audit reinforces this recommendation rather than complicating it:
because the most accurate method on the median, the per-scene SVM, is also
the most expensive (Section~\ref{sec:res_cost}), its added complexity buys
no general gain, and the simple per-scene VV Otsu is the more practical
default.

\subsection{A/P as an a-priori screening tool}

Because A/P is computed from a water-body polygon already available in global
databases (GRanD, GRDL, or the JRC extent), it can be evaluated in advance of
processing. Its value is that it predicts the geometric ceiling on accuracy:
high-A/P reservoirs are tracked reliably, whereas low-A/P reservoirs are only
permitted, not guaranteed, to fail (Section~\ref{sec:res_ap}). A water manager
can therefore rank a portfolio of reservoirs by expected SAR reliability before
acquiring or processing any imagery, and direct field effort or complementary
altimetry to the low-A/P systems where errors are likely. The screening is
inexpensive to deploy, since the training samples and the reservoir boundary are
drawn automatically from global water and land-cover masks, which removes the
manual delineation that would otherwise limit transfer to new regions. Coupled
with the parsimony result above, a single simple detector applied everywhere,
with A/P used only to flag where its output should be trusted, keeps the
operational workflow as light as possible.

This a-priori value is not incidental to A/P's particular form. Other
shoreline-shape descriptors, such as fractal dimension or compactness
indices, also quantify shoreline irregularity, but they are not derived
from the mixed-pixel mechanism itself (Section~\ref{sec:introduction}), so
there is no equivalent argument for why they, rather than some other
irregularity measure, should set an accuracy ceiling. A/P is instead the
length-scale that mechanism directly implies: the ambiguous-boundary
fraction $r/(A/P)$ follows from dividing the boundary-strip area by the
water area, so the relationship it predicts, and that we observe (Spearman
$\rho=0.51$, Section~\ref{sec:res_ap}), is consistent with that derivation
rather than a post-hoc fit to whichever shape metric happened to correlate.
We did not screen candidate shoreline descriptors and keep the
best-performing one; A/P was derived from the error mechanism first, and
only then tested.

\subsection{Geometric and radiometric limits}
\label{sec:disc_geo_radio}

A/P bounds accuracy from geometry and is blind to radiometric effects, which
form a second axis. The clearest radiometric candidate, wind, does not organise
along geometry and shows no systematic effect on the 47-reservoir set tested
(Section~\ref{sec:res_climate}). One reason the Bragg-scattering prediction
\citep{Valenzuela1978} may not surface here is that Otsu's threshold is
itself re-estimated every scene (Section~\ref{sec:methods_otsu}): a decision
rule that recalibrates each acquisition can absorb some of the wind-driven
backscatter shift before it becomes an omission error, in a way that Bragg's
original formulation, developed against a fixed detection threshold, does not
anticipate. This is a sample-specific finding, not a general one, and a
larger or more wind-exposed set could still reveal an effect the present test
is not powered to detect. The intuitive attribution of specific
high-A/P failures to wind or ice does not hold up, and one such case, Rockport,
is in fact well tracked by both per-scene detectors once the arbitrary
fixed-training baseline is set aside. The residual that A/P does not explain is
instead organised, in part, by climate: reservoirs in humid, vegetated settings
are tracked less well at given A/P than those in arid or Mediterranean ones.
Reaching this conclusion required first separating a genuine climate signal
from a data-quality artefact. In a first pass, the climate residual lost
significance once reservoirs with a demonstrably noisy optical reference were
excluded after the fact, since those reservoirs were disproportionately
humid and tropical, and we could not tell whether climate mattered or the
reference did. Only once the reference-quality screen (Section~\ref{sec:val_global})
was applied throughout, rather than as a post-hoc exclusion, did the
climate residual prove robust: on this cleaned set it remains
significant ($p=0.012$), so it is not merely standing in for reference noise:
climate is a genuine, if secondary, modulator of the geometric ceiling A/P
sets. The practical implication survives either way. Where two independent
SAR detectors agree on a stable reservoir and an optical reference does not,
the disagreement is a limit of the reference, not of the SAR, and this is
precisely where an all-weather SAR record is most valuable.

\subsection{Limitations}
\label{sec:disc_limitations}

Several limitations bound these conclusions. The study addresses surface area
only. Converting area to storage volume requires an area--elevation--volume
relationship that, for a significant number of reservoirs, needs updating
or recreating from scratch (Section~\ref{sec:introduction}); we leave that
conversion to future work.
The global set is referenced
against the 30~m JRC product, an inter-product comparison rather than ground
truth. The four Sicilian reservoirs validated against 3~m PlanetScope reproduce
the main findings, but a broader near-truth validation would strengthen them,
especially in the humid settings where the optical reference is weakest. Four
reservoirs whose regulated level is nearly constant were excluded a priori,
since a flat reference makes KGE undefined in practice. They are monitored
perfectly well but cannot be scored this way. The reference-quality screen
itself, while principled, uses a single threshold (roughness ratio
$\geq$2.5) calibrated on this dataset, and a stricter or looser cut would move a
handful of borderline reservoirs across the line.

Motivated by the Ancipa mechanism (Section~\ref{sec:disc_geo_radio}), we also
tested whether A/P computed dynamically, per scene, from the SAR mask itself
could serve as a live operational trigger for switching from Otsu to the SVM.
A genuine, narrow effect exists: in a $\sim$51--70~m dynamic-A/P window, the
SVM's monthly area error is typically smaller than Otsu's, sharpening the
mixed-pixel mechanism at the most extreme drawdowns. It does not, however,
translate into a usable rule. The effect does not hold within individual
reservoirs' own trajectories (the median within-reservoir correlation between
dynamic A/P and the SVM's monthly edge is slightly positive, not negative, as
the mechanism would require), and a switching policy backtested at that
threshold more often lowers the resulting KGE than raises it, because a
monthly series assembled from two differently-biased detectors introduces
discontinuities that a correlation-based metric penalises heavily, even where
the segment the SVM contributes is individually more accurate. We could not
identify any variable, geometric or otherwise (static or dynamic A/P, climate,
or the detector's own baseline skill), that predicts in advance which
reservoirs benefit most from the SVM, so a per-scene or per-reservoir
switching rule is not yet justified. This leaves the roughly 40\% of
reservoirs where the SVM's advantage is large (Section~\ref{sec:res_fourway})
without a confirmed explanation. Untested candidates include riparian
vegetation, to which the cross-polarised channel is differentially sensitive,
shoreline complexity beyond the single A/P scalar (for example the number of
disconnected sub-basins or arms at a given A/P), seasonal ice or snow, and
per-reservoir variation in training-sample quality. Finding that variable, or
a splicing-aware combination rule that does not penalise the transition
itself, is left to future work.

Finally, the temperate/continental climate zone's low median is not fully
explained. Several of its worst-performing members (including one hydropeaking
reservoir with a documented SVM-training phase lag, and two subsequently found
to fail the reference-quality screen) have identifiable, non-climatic causes,
and removing every such case still leaves a modest, only marginally significant
residual driven largely by narrow, valley-confined hydropower reservoirs in
central Europe. These reservoirs share regulated, operationally-driven water
levels as much as they share a climate zone, so an operational confound
(active drawdown-refill cycling for hydropower generation, which can decouple
short-term level changes from the seasonal signal A/P-based expectations
implicitly assume) is at least as plausible as a radiometric climate mechanism.
We report the climate effect honestly as significant on the cleaned data,
while flagging that its temperate-zone component may reflect reservoir
operation rather than climate per~se. Distinguishing the two is left to future
work with an explicit operational-regime covariate.
